\begin{afterword}
\summaryhead

The author recalls the moment, at fifteen, when the journey as a rationalist began. The fifteen-year-old was going over a pleasing memory from summer camp at six or seven. A teenage counselor had proposed a game in which the boys spanked each other, and the author had refused and been sent to sit in the corner. The memory had come to stand for an early refusal to hurt others and to play a negative-sum game.

At fifteen the author saw that this was untrue. A seven-year-old knew nothing of the Prisoner's Dilemma. The real reason for refusing was not wanting to be hurt, and the author had always known it, glancing at the real memory and looking away.

Having caught that feeling, the author resolved to notice every time an unwanted truth was pushed aside, and named the skill singlethink, after Orwell's doublethink: notice that you are forgetting, and remember. The author doubts it is original. Confirmation bias runs deep, the post ends, but that first effort cleared many corners.

\responsehead

The post is a short memoir that names a skill. It gives its case in full, it is candid about the vagueness of the childhood memory, and it concedes that the skill may not be new. Its gloss on Orwell is accurate. The problems are in the step that corrects the memory, and in the claim that the skill worked.

The correction has less support than it seems. The flattering reading had two parts: the child would not hurt others, and the child refused a negative-sum game. The post's one piece of evidence, that a seven-year-old could not have known the Prisoner's Dilemma, disposes of the second part only. A child can refuse to hurt others without any game theory. The choice between that motive and ``I didn't want to get hurt'' is made by introspection: the fifteen-year-old's sense of having ``always known.''

The post has just shown that the author's own reading of this memory could be confidently wrong. It does not ask whether the new reading could be wrong too, and it gives no way to tell a truth recovered in this way from a new story that simply feels true. The less flattering reading may well be the right one. The post does not show it.

The skill has the same problem. It depends on a feeling, the feeling of shoving an unwanted truth into a corner. The post does not say how often that feeling is mistaken, or how one would find out. Its only evidence of success is the closing sentence: the first ``broom'' swept ``quite a few corners.'' No second corner is described.

In short: a candid memory that names a skill, with its central correction settled by a feeling of having always known, and its success reported without a second case.
\end{afterword}
